\documentclass[10pt]{article}
\usepackage[T1]{fontenc}
\usepackage[utf8]{inputenc}
\usepackage[french]{babel}
\usepackage{amsmath,amssymb,mathrsfs,mathpazo}
\usepackage{geometry}
\geometry{a4paper,margin=2cm}
\usepackage{graphicx,array,tabularx,tabularray,colortbl,multirow,multicol,booktabs,enumitem,pifont,ulem,fancybox,diagbox,cancel}
\usepackage{pstricks,pstricks-add,pst-plot,pst-tree,pst-node,pst-text,pst-func,pst-math,pst-bspline,pst-3dplot,pst-eucl,pst-solides3d}
\let\pscalculate\undefined
\usepackage{pst-bezier}
\usepackage{tikz,pgfplots,tkz-tab}
\usetikzlibrary{arrows,arrows.meta,patterns,calc,positioning,trees,decorations.markings,shapes.geometric,shapes.misc,shapes.arrows,intersections,decorations.pathmorphing,decorations.pathreplacing,backgrounds,shadows,babel}
\pgfplotsset{compat=1.18}
\usepgfplotslibrary{fillbetween}
\usepackage[np]{numprint}
\usepackage{listings,makecell,pas-tableur,twoopt}
\let\barre\undefined
\usepackage[pstricks]{tablvar}
\definecolor{darkgreen}{RGB}{0,100,0}
\providecommand{\red}{\color{red}}
\providecommand{\blue}{\color{blue}}
\providecommand{\green}{\color{green}}
\providecommand{\black}{\color{black}}
\providecommand{\white}{\color{white}}
\providecommand{\gray}{\color{gray}}
\providecommand{\cyan}{\color{cyan}}
\providecommand{\magenta}{\color{magenta}}
\providecommand{\yellow}{\color{yellow}}
\providecommand{\violet}{\color{violet}}
\providecommand{\orange}{\color{orange}}
\providecommand{\euro}{\texteuro}
\providecommand{\EUR}{\texteuro}
\providecommand{\ieme}{\textsuperscript{e}}
\providecommand{\degres}{\ensuremath{^\circ}}
\providecommand{\degree}{\ensuremath{^\circ}}
\providecommand{\vv}[1]{\overrightarrow{#1}}
\providecommand{\Coord}[1]{#1}
\providecommand{\decosix}{}
\providecommand{\starredbullet}{\ensuremath{\bullet}}
\pagestyle{empty}
\setlength{\parindent}{0pt}
\usepackage{adjustbox}\DeclareUnicodeCharacter{2236}{:}\DeclareUnicodeCharacter{2190}{\ensuremath{\leftarrow}}
\providecommand{\R}{\mathbb{R}}
\providecommand{\N}{\mathbb{N}}
\providecommand{\Z}{\mathbb{Z}}
\providecommand{\Q}{\mathbb{Q}}
\providecommand{\D}{\mathbb{D}}
\providecommand{\C}{\mathbb{C}}
\providecommand{\e}{{\mathrm{e}}}
\providecommand{\vect}[1]{\overrightarrow{#1}}
\providecommand{\vectt}[1]{\overrightarrow{\mathrm{#1}}}
\providecommand{\vvt}[1]{\overrightarrow{\mathrm{#1}}}
\providecommand{\np}[2][]{{#2~#1}}
\providecommand{\vv}[1]{{\overrightarrow{#1}}}
\renewcommand{\d}{{\mathrm{d}}}
\providecommand{\degres}{{^\circ}}
\providecommand{\degre}{{^\circ}}
\providecommand{\blue}{\color{blue}}
\providecommand{\red}{\color{red}}
\providecommand{\enskip}{\quad}
\providecommand{\llbracket}{\lbrack\!\lbrack}
\providecommand{\rrbracket}{\rbrack\!\rbrack}
\providecommand{\euro}{{\text{€}}}
\providecommand{\no}{{\text{n°}}}
\providecommand{\interval}[2]{{[#1;#2]}}
\providecommand{\intervFO}[2]{{[#1;#2[}}
\providecommand{\dotfill}{\ldots}

\providecommand{\ds}{\displaystyle}
\providecommand{\textding}[1]{#1}
\providecommand{\Oij}{\(\mathrm{O};\vec{\imath},\vec{\jmath}\)}
\providecommand{\Oijk}{\(\mathrm{O};\vec{\imath},\vec{\jmath},\vec{k}\)}
\providecommand{\Ouv}{\(\mathrm{O};\vec{u},\vec{v}\)}
\providecommand{\intervFF}[2]{\ensuremath{[#1;#2]}}
\renewcommand{\euro}{\ifmmode\text{\texteuro}\else\texteuro\fi}
\DeclareUnicodeCharacter{00B0}{\ifmmode^{\circ}\else\textdegree\fi}
\DeclareUnicodeCharacter{20AC}{\ifmmode\text{\texteuro}\else\texteuro\fi}
% Énoncés : Éducation nationale. Transcription et corrigés : APMEP, Denis Vergès et contributeurs. Découpage et classement : CoopMaths, Sébastien Lozano, Éric Elter et contributeurs. Adaptation de la présentation : Maths Explorer.
% CC BY-SA : https://coopmaths.fr/www/about/
% https://coopmaths.fr/alea/static/e3c/2020/tex/e3c_2020_00_sujet8_3_cor.tex
\begin{document}
\noindent\textit{APMEP / CoopMaths. CC BY-SA. Présentation adaptée par Maths Explorer.}\par\medskip
\section*{Exercice 3 (5 points)}
	
	
	\subsection*{1. En exécutant le script Python ci-dessous, on obtient la liste [0, 1, 1, 1, 1, 1, 0].}
	
Le médicament est efficace de la première à la cinquième heure.
	
	\subsection*{2. On admet que la fonction $f$ est dérivable sur l’intervalle $[0; 6]$, calculer sa fonction dérivée.}
	
La fonction polynôme $f$ est dérivable sur $\mathbb{R}$, donc sur $[0; 6]$ :
	\[
	f'(x) = 3x^2 - 24x + 36 = 3(x^2 - 8x + 12)
	\]
	
	\subsection*{3. Justifier que la tangente $T$ à la courbe représentative de la fonction $f$ au point $A$ d’abscisse 4 admet pour équation réduite $y = -12x + 64$.}
	
On a $f(4) = 16$ et $f'(4) = -12$.\\
Une équation de la tangente au point de coordonnées $(4, 16)$ est donc :
	\[
	y - 16 = -12(x - 4) \quad \text{ou} \quad y = 16 - 12x + 48 \quad \text{et finalement} \quad y = -12x + 64
	\]
	
	\subsection*{4. Démontrer que $f(x) - (-12x + 64) = (x - 4)^3$.}
	
	On a :
	\[
	f(x) - (-12x + 64) = x^3 - 12x^2 + 36x - (-12x + 64) = x^3 - 12x^2 + 36x + 12x - 64
	\]
	
	Or
	\[
	(x - 4)^3 = (x - 4)^2 (x - 4) = (x^2 + 16 - 8x) (x - 4) = x^3 - 4x^2 + 16x - 64 - 8x^2 + 32x = x^3 - 12x^2 + 48x - 64
	\]
	
	Donc
	\[
	f(x) - (-12x + 64) = (x - 4)^3
	\]
	
	\subsection*{5. En déduire la position relative de la courbe représentative de la fonction $f$ par rapport à la tangente $T$ au point $A$.}
	
	Comme $(x - 4)^3 = (x - 4)^2 (x - 4)$ et que $(x - 4)^2 > 0$ quel que soit le réel $x$, le signe de $(x - 4)^3$ est celui de $(x - 4)$. Donc
	\begin{itemize}
		\item Sur $[0; 4[$, $x - 4 < 0$, donc $f(x) - (-12x + 64) < 0$ ce qui signifie géométriquement que la courbe est en dessous de sa tangente ;
		\item Sur $]4; 6[$, $x - 4 > 0$, donc $f(x) - (-12x + 64) > 0$ ce qui signifie géométriquement que la courbe est au-dessus de sa tangente ;
		\item Pour $x = 4$, la courbe et sa tangente ont la même ordonnée.
	\end{itemize}

\end{document}
